% ---------------------------------------------------------------- remerciements
\chapter*{Remerciements}
\addcontentsline{toc}{chapter}{Remerciements}

Je tiens à exprimer ma profonde gratitude à \textbf{\encadrant}, enseignant au Département de
Mathématiques et Informatique de la Faculté des Sciences et Techniques, pour la proposition de ce
sujet, la qualité de son encadrement et la pertinence de ses orientations tout au long de ce projet.

Mes remerciements s'adressent également à l'ensemble du corps professoral de la section Informatique
de l'Université Cheikh Anta Diop de Dakar, et particulièrement aux enseignants du Master II option
Business Intelligence, pour les connaissances transmises durant ce cursus.

Je remercie enfin ma famille et mes camarades de promotion pour leur soutien et leurs encouragements.

% ---------------------------------------------------------------- résumé
\chapter*{Résumé}
\addcontentsline{toc}{chapter}{Résumé}

Les organisations disposent de volumes croissants de données opérationnelles, mais leur exploitation
reste freinée par la complexité technique des outils d'interrogation : sans maîtrise du langage SQL, les
décideurs dépendent des équipes techniques pour obtenir le moindre indicateur.

Ce projet met en œuvre une \textbf{Modern Data Platform} complète sur le jeu de données de référence
\emph{Wide World Importers} de Microsoft. Les données transactionnelles de \textbf{SQL Server} sont
ingérées de manière incrémentale dans \textbf{Snowflake} par \textbf{DLT}, puis organisées selon une
\textbf{architecture Médaillon} : une couche \emph{Bronze} brute (30 tables), une couche \emph{Silver}
nettoyée (29 modèles) et une couche \emph{Gold} modélisée en étoile (15 tables, dont trois historisées
en SCD de type 2). Les transformations sont développées, documentées et testées avec \textbf{dbt}
(189 nœuds, dont plus de 140 tests de qualité), et l'ensemble est orchestré par \textbf{Dagster}
(75 assets, 146 contrôles, exécution quotidienne avec reprise sur erreur et alertes).

Une application \textbf{Streamlit} permet enfin aux décideurs d'interroger la couche Gold
\textbf{en langage naturel} : les questions, posées en français, sont traduites en SQL par
\textbf{Snowflake Cortex Analyst} grâce à un modèle sémantique décrivant les tables, leurs relations,
leurs synonymes métier et leurs indicateurs. La requête générée est contrôlée (lecture seule), exécutée
avec un rôle aux droits restreints, puis restituée sous forme de tableau et de graphique. L'application
est déployée à la fois en local et dans \emph{Streamlit in Snowflake}, et une chaîne d'intégration
continue GitHub Actions valide chaque modification sur un clone « zéro copie » de l'entrepôt.

\textbf{Mots-clés :} Modern Data Stack, architecture Médaillon, Snowflake, DLT, dbt, Dagster,
modélisation dimensionnelle, SCD2, Text-to-SQL, Cortex Analyst, LLM, Streamlit.

% ---------------------------------------------------------------- abstract
\chapter*{Abstract}
\addcontentsline{toc}{chapter}{Abstract}

\begin{otherlanguage}{english}
Organisations accumulate ever-growing volumes of operational data, yet business users who do not master
SQL still depend on technical teams to obtain even simple indicators.

This project builds an end-to-end \textbf{Modern Data Platform} on Microsoft's \emph{Wide World
Importers} sample database. Transactional data from \textbf{SQL Server} is incrementally ingested into
\textbf{Snowflake} with \textbf{DLT} and organised following a \textbf{Medallion architecture}
(Bronze, Silver and a Gold star schema including type-2 slowly changing dimensions). Transformations
are versioned, documented and tested with \textbf{dbt}, and the whole pipeline is orchestrated by
\textbf{Dagster} with daily scheduling, retries, freshness checks and failure alerts.

A \textbf{Streamlit} application lets decision makers query the Gold layer in \textbf{natural
language}: French questions are translated into SQL by \textbf{Snowflake Cortex Analyst} using a
semantic model, executed under a read-only role, and rendered as tables and charts. The application runs
locally and in \emph{Streamlit in Snowflake}, and a GitHub Actions pipeline validates every change on a
zero-copy clone of the warehouse.

\textbf{Keywords:} Modern Data Stack, Medallion architecture, Snowflake, DLT, dbt, Dagster,
dimensional modelling, SCD2, Text-to-SQL, Cortex Analyst, LLM, Streamlit.
\end{otherlanguage}

% ---------------------------------------------------------------- sigles
\chapter*{Liste des sigles et abréviations}
\addcontentsline{toc}{chapter}{Liste des sigles et abréviations}

\begin{longtable}{@{}p{3cm}p{11.5cm}@{}}
\toprule
\textbf{Sigle} & \textbf{Signification} \\
\midrule
\endhead
BI & Business Intelligence (informatique décisionnelle) \\
CA & Chiffre d'affaires \\
CI/CD & Intégration continue / Déploiement continu \\
DAG & Directed Acyclic Graph (graphe orienté acyclique) \\
dbt & data build tool \\
DLT & data load tool \\
ELT & Extract, Load, Transform \\
ETL & Extract, Transform, Load \\
HT / TTC & Hors taxes / Toutes taxes comprises \\
KPI & Key Performance Indicator (indicateur clé de performance) \\
LLM & Large Language Model (grand modèle de langage) \\
MDP & Modern Data Platform \\
OLAP & Online Analytical Processing \\
OLTP & Online Transaction Processing \\
RBAC & Role-Based Access Control \\
SCD & Slowly Changing Dimension (dimension à évolution lente) \\
SiS & Streamlit in Snowflake \\
SQL & Structured Query Language \\
TVA & Taxe sur la valeur ajoutée \\
WWI & Wide World Importers \\
YAML & YAML Ain't Markup Language \\
\bottomrule
\end{longtable}
